\chapter{ENVIRONMENTAL STUDY}

\hspace{0.5cm}
This chapter outlines the hardware and software configuration environments required to successfully implement, deploy, and execute the RustAuditAI system. It details the platform constraints, compiler front-end parser utilities, language inference tools, and graph visualization libraries that form the underlying stack of the multi-dimensional software quality intelligence pipeline.

\section{System Configuration}
\hspace{0.5cm}
The operational environment for RustAuditAI spans specific physical computing architectures and virtual software execution layers designed to handle high-fidelity abstract syntax parsing and local model inference workflows efficiently.

\subsection{Hardware Requirements}
\hspace{0.5cm}
The hardware requirements outline the minimum resources necessary to operate the system. These resources are crucial for performing code analysis, processing tasks and running machine learning operations.

The key hardware requirements are detailed as follows:


\begin{itemize}
\item \textbf{Processor}: Intel\textsuperscript{\textregistered} Core\textsuperscript{\texttrademark} i5 Processor / AMD Ryzen\textsuperscript{\texttrademark} 5 Processor or higher, to efficiently handle code processing tasks.

\item \textbf{Memory}: Minimum 8 GB RAM, to ensure smooth execution of analysis and model operations.

\item \textbf{Storage}: 50 GB of free disk space, for storing code files and generated outputs.
\end{itemize}

\subsection{Software Requirements}
\hspace{0.5cm}
The requirements for software specify the necessary platforms, languages and libraries that will be used in running the system. The use of these tools, activities like parsing, integration and documentation become easy. Compatibility assurance is relevant to ensure that the various components of the system communicate effectively. 


The key software requirements are detailed as follows:
\begin{itemize}
    \item \textbf{Operating System}: Linux distribution (Ubuntu 22.04 LTS or later recommended), macOS Sequoia, or Microsoft Windows 11 configured with Windows Subsystem for Linux (WSL2).
    
    \item \textbf{Primary Implementation Platforms}: Python runtime environment version 3.11.x or later, alongside the native Rust compiler toolchain (rustc and cargo engine package utility version $\ge$ 1.75).
    
    \item \textbf{Version Control System}: Git open-source distributed source pipeline control suite.
    
    \item \textbf{Artificial Intelligence Interface Model}: Access token credentials for the centralized Groq inference platform API cluster utilizing the LLaMA-3.3-70B instruction-tuned model parameter footprint.
\end{itemize}

\subsection{Software Requirements for Development}
\hspace{0.5cm}Besides the execution specifications, some indispensable are the tools for development, testing and maintenance. They assist in writing and debugging programs, as well as managing the code appropriately. Besides, the tools help in maintaining versions and assembling different systems together. 

The key development software requirements for development are detailed as follows:
\begin{itemize}
    \item \textbf{Operating System}: Windows 10, Ubuntu 20.04 LTS or macOS, ensuring cross-platform compatibility.
    
    \item \textbf{Back End}: Python 3.8 or later, serving as the core implementation language.
    
    \item \textbf{IDE Used}: Visual Studio Code, Antigravity IDE used for development and debugging tasks.
    
    \item \textbf{Version Control}: Git open-source distributed source pipeline control suite.
    
    \item \textbf{Development Tools and Frameworks}:
    \begin{itemize}
        \item \textbf{syn crate:} Specialized parsing framework utilized to read and decode token streams into explicit Rust abstract syntax trees (ASTs).
        \item \textbf{quote and proc-macro2 crates:} To transform structured AST tokens back into compiler-valid syntactic elements for analysis slicing.
        \item \textbf{Pydantic crate/library:} Tasked with parsing and structurally enforcing input/output JSON schemas coming from the text reasoning layers.
        \item \textbf{NetworkX library suite:} Deployed as the core runtime graph data structure layout utility to handle localized nodes, edges, and path traversals for the Function-Level Ownership Graph (FLOG).
        \item \textbf{Requests or Httpx engine packages:} Configured to manage safe, asynchronous connectivity protocols with cloud API endpoints.
        \item \textbf{Rich text layout library:} Configured to parse code diagnostics into interactive terminal layout screens.
    \end{itemize}
\end{itemize}
\section{Software Specification}
\hspace{0.5cm}The following section provides a detailed description of different software components and technologies that are used to design the system under discussion. This is because each of the software components performs specific tasks that shall be useful when analyzing the code and preparing the documentation. Each software component is made to perform particular operations that are useful in code analysis and documentation preparation and the use of technologies is meant to ensure a smooth integration process. 

\subsection{Python 3.8}
\hspace{0.5cm}
Python functions as the primary execution coordinator for RustAuditAI. It orchestrates the lifecycle of the static analysis pipeline, handling file validation, initiating external compilation checks, managing the graph data structures inside NetworkX, and routing extracted semantic data vectors to the remote inference gateways. Python’s mature data structure packages and rich reporting wrappers make it the ideal platform to bind the deterministic compiler graph layers with the downstream explainable text engines.

\subsection{Rust Toolchain(syn and cargo)}
\hspace{0.5cm}
The native Rust ecosystem acts as the foundational parsing source engine for this framework. The syn crate acts as a robust compiler front-end parser that ingests standard .rs file text strings and decodes them into fully structural Abstract Syntax Trees (ASTs). This component preserves localized code expressions, statement mappings, control constructs, and identifier paths. By relying on syn, RustAuditAI gains direct access to precise compiler-level representations of function scopes without requiring fully compiled intermediate binary builds, establishing the raw data foundation for the custom Ownership and Control Flow graphs.

\subsection{The Function-Level Ownership Graph (FLOG) Engine}
\hspace{0.5cm}
The FLOG Engine is a core domain-specific novelty created for this project using the NetworkX library. It translates abstract code definitions into explicit variable lifecycle tracking paths. It generates nodes representing allocation footprints (such as local references, box vectors, and variable re-bindings) interconnected by edges detailing ownership rules (such as moves, mutable borrows, or immutable aliases). This custom data structure enables the static analysis engine to perform deep intra-procedural lifecycle evaluations, mathematically identifying code flaws like allocation leaks or excessive cloning behaviors before execution.

\subsection{Git}
\hspace{0.5cm}
Git is used as the version control system for analyzing the structure of the repository and obtaining historical information regarding the code. This allows the system to work with entire projects by cloning repositories from remote sites, which eliminates manual uploading of separate files. Information at the commit level is also available, for example, for the production of documentation, such as a list of the latest changes. 

\par This would enable the system to be able to identify the different types of changes through commits, thus being able to identify development patterns and get knowledge on how the code is being changed over time. All this would lead to context being better understood, thereby adding a layer of value to the kind of document that is being generated. Other than source code repositories, the system would also allow local source code directories to ensure flexibility with the various types of input sources. 

\subsection{NetworkX}
\hspace{0.5cm}
It is used to create and analyze the graphs representing the relationships, function calls and dependencies in the code. NetworkX helps the system to model the code structure as a directed graph, where the functions are represented as nodes and the interactions are represented as edges. This allows the system to understand important parts of the program code, such as the starting point, recursion and any dependencies. The use of graph-based analysis increases the understanding of how the code operates. Additionally, the system uses graph analysis to perform complex evaluations of the code, enabling a better understanding of its critical components.

\subsection{Explainable AI (XAI) Module}
\hspace{0.5cm}
The Explainable AI module improves operational transparency by turning raw, numeric quality checks into clear, natural language recommendations. Rather than printing standard compiler diagnostics or raw text warning codes, this component uses structured semantic properties to explain why an expression violates best practices. It explicitly documents the data path where a variable bottleneck emerges and includes structural feedback explaining how a suggested update modifies execution dependencies. This detailed explanation turns abstract graph diagnostics into educational, highly actionable refactoring advice for developers.

\subsection{Visualization Support}
\hspace{0.5cm}
The visualization engine translates internal graph structures into self-contained HTML reports featuring interactive Control Flow Graphs (CFGs) and variable lifecycle grids. Built without external layout software dependencies, it provides custom visual highlights along vulnerable or inefficient execution pathways inside a function body. By showing a side-by-side behavioral comparison between original and proposed code structures, the module enables developers to trace attack vectors or performance bottlenecks visually, improving system transparency and patch management.